\documentclass[11pt]{article}
\usepackage[T1]{fontenc}
\usepackage[utf8]{inputenc}
\usepackage{lmodern}
\usepackage[margin=1in]{geometry}
\usepackage{amsmath,amssymb,amsthm}
\usepackage[title]{appendix}
\usepackage[colorlinks=true,linkcolor=blue,citecolor=blue,urlcolor=blue]{hyperref}
\hypersetup{pdfauthor={P. Graham},pdftitle={Connected quandles that refute Hayashi's conjecture}}
\newtheorem{theorem}{Theorem}[section]
\newtheorem{lemma}[theorem]{Lemma}
\newtheorem{proposition}[theorem]{Proposition}
\newtheorem{corollary}[theorem]{Corollary}
\newtheorem*{hconj}{Hayashi's conjecture}
\newcommand{\F}{\mathbb{F}}
\title{Connected quandles that refute Hayashi's conjecture}
\author{P. Graham\\ \small\url{https://orcid.org/0009-0007-8493-1505}}
\date{November 5, 2026}
\begin{document}
\maketitle
\begin{abstract}
We construct connected quandles of orders 1890 and 5670 whose right
translations $x\mapsto x*y$ have cycles of lengths 60, 84, 140 and 210.
Since 60, 84 and 140 do not divide 210, this disproves the conjecture of
Hayashi (2013) that in a finite connected quandle the longest cycle length
of a right translation is a multiple of all the others. A quandle is a set
with a self-distributive operation modeled on conjugation in a group, and
quandles give invariants of knots. The conjecture was known for profiles
with at most five cycles (Lages and Lopes), for orders up to 47 (Watanabe)
and for affine quandles (Kajiwara and Nakayama). Filipi showed that for
conjugation quandles it is equivalent to a group-theoretic conjecture,
that a power of an element of a finite generating conjugacy class $C$ is
central as soon as it commutes with a suitable element of $C$. Our
quandles are conjugacy classes of affine maps, so they refute that
conjecture too. Each cycle length is twice the product of only three of
the primes 2, 3, 5 and 7, while each right translation has order 420. The
construction gives an infinite family, including latin counterexamples
(every left translation is also a bijection), and no counterexample from
it has order below 1890. The proofs are by hand. Independent programs
recompute both quandles, and one checks every quandle axiom on the full
operation table.

This work has not yet been reviewed by experts. The proofs, the checking programs and the text were produced by AI models (from OpenAI and Anthropic), working under the author's direction.
\end{abstract}
\medskip
\noindent\textbf{MSC 2020:} 20N02, 57K12, 20D45.\quad
\textbf{Keywords:} quandle, profile, Hayashi's conjecture, conjugacy class, cycle type.

\section{Introduction}

A \emph{quandle} is a set $Q$ with a binary operation $*$ such that
$x*x=x$, every right translation $R_y\colon x\mapsto x*y$ is a bijection,
and $(x*y)*z=(x*z)*(y*z)$ for all $x,y,z\in Q$. It is \emph{connected} if
the group generated by the right translations acts transitively on $Q$,
\emph{faithful} if $y\mapsto R_y$ is injective, and \emph{latin} if every
left translation $x\mapsto y*x$ is a bijection. The basic example is a
subset $C$ of a group that is closed under conjugation by its elements and
their inverses, with $x*y=yxy^{-1}$; then $R_y$ is conjugation by $y$.
This is a \emph{conjugation quandle}.

Since $R_{x*y}=R_yR_xR_y^{-1}$, in a connected quandle all right
translations are conjugate, so they share a cycle type. Following Lages
and Lopes~\cite{LL}, the \emph{profile} of a finite connected quandle is
this cycle type, written $(\ell_1,\dots,\ell_r)$ with
$1=\ell_1\le\dots\le\ell_r$ and repeats allowed; we also write $a^m$ for
$m$ cycles of length $a$, so that $(1,2,2,6)$ is $1\,2^2\,6$. Hayashi
posed the following as \cite[Conjecture~1.1]{Ha} (arXiv 2011, published
2013); we state it as in \cite[Conjecture~1.6]{LL}, with $r$ for the number of cycles.

\begin{hconj}
Let $(X,*)$ be a connected quandle of order $n$ with profile
$(\ell_1,\dots,\ell_r)$, where $1=\ell_1\le\dots\le\ell_r$. Then $\ell_r$
is a multiple of $\ell_i$ for every $i\in\{1,\dots,r\}$.
\end{hconj}

\subsection{Main results}

\begin{theorem}\label{thm:main}
There are connected, faithful quandles of orders $1890$ and $5670$, with
profiles
\[
1^{2}\,2^{104}\,60^{7}\,84^{5}\,140^{3}\,210^{2}
\qquad\text{and}\qquad
1^{2}\,2^{314}\,60^{21}\,84^{15}\,140^{9}\,210^{6},
\]
respectively. Since $60$, $84$ and $140$ do not divide $210$, Hayashi's
conjecture is false. Both quandles are conjugation quandles: each is a
single finite conjugacy class of the group it generates.
\end{theorem}

The two quandles are given in Section~\ref{sec:examples} and
Appendix~\ref{app:data}. Their profiles have $123$ and $367$ cycles, and
both orders exceed $47$. Both meet the necessary condition of
\cite[Corollary~2.6]{LL} for a counterexample: $210$ divides
$\operatorname{lcm}(60,84,140)=420$. The quandle of order $1890$ is the
smallest the construction can give (Proposition~\ref{prop:min}), and it has
the profile of the case $p=3$, $\{q_j\}=\{2,3,5,7\}$ of
Theorem~\ref{thm:family} below. The quandle of order $5670$ differs only at
the prime $3$, where it uses the regular module $\F_3[x]/(x^3-1)$; it shows
that the module changes the order but not the cycle lengths.

Filipi~\cite{F} calls a finite conjugacy class $C$ of a group \emph{good}
when every $c\in C$ admits some $z\in C$ with
$\langle c\rangle\cap C_G(z)\le Z(\langle C\rangle)$
\cite[Definition~4.1]{F}. Filipi proved that Hayashi's conjecture holds for
all connected conjugation quandles if and only if every finite conjugacy
class $C$ with $\langle C\rangle=G$ is good \cite[Theorem~4.1]{F}, a work
that ``completes Kayacan's ideas that first appeared in an arXiv paper [8]
in 2021'' \cite[\S1]{F}, with reference to \cite{K}. The condition is
stated as \cite[Conjecture~4.1]{F}: ``Let $G$ be a group and $C$ its
finite conjugacy class such that $\langle C\rangle=G$. Then for any $c\in
C$, there exists some $z\in C$ such that for all $k\in\mathbb{Z}$, the
element $c^k$ is central if and only if it commutes with $z$.''

\begin{corollary}\label{cor:filipi}
Filipi's Conjecture~4.1 is false: for each quandle $C$ of
Theorem~\ref{thm:main} and $H=\langle C\rangle$, the conclusion fails at
every $c\in C$.
\end{corollary}

The construction uses the group $\F_3^2$, whose four subgroups of order
$3$ each switch off one of the primes $2$, $3$, $5$, $7$ in the cycle
lengths. Replacing it by $\F_p^2$, with its $p+1$ subgroups of order $p$
and $p+1$ primes, gives:

\begin{theorem}\label{thm:family}
Let $p$ be an odd prime and $q_0,\dots,q_p$ distinct primes. There is a
connected, faithful quandle of order $p^{2+\varepsilon}\prod_{q_j\ne p}
q_j^{(p-1)/2}$, where $\varepsilon=1$ if $p\in\{q_0,\dots,q_p\}$ and
$\varepsilon=0$ otherwise, whose profile consists of the lengths $1$, $2$
and $\ell_j=2\prod_{i\ne j}q_i$ for $0\le j\le p$. Exactly $p$ of the $p+1$
distinct lengths $\ell_j$ fail to divide the longest one. If all $q_j$
are odd, the quandle is latin.
\end{theorem}

So the number of cycle lengths that fail to divide the longest one is
unbounded. The smallest latin case, $p=3$ and $\{q_j\}=\{3,5,7,11\}$, has
order $10395$ and profile $1\,2^{577}\,210^{11}\,330^{7}\,462^{5}\,770^{3}$.

\begin{proposition}\label{prop:min}
Let $C$ be the conjugation quandle of Section~\ref{sec:setting}, built
from any group $A$ of odd order acting on a sum of vector spaces over
distinct prime fields, any involution $P$ with $PaP=a^{-1}$ for $a\in A$,
and any $g_0$. If some cycle length of $C$ does not divide the longest,
then $|C|\ge1890$, and $|C|\ge10395$ if $C$ is latin. Both bounds are
attained.
\end{proposition}

The proof is in Appendix~\ref{app:min}.

\subsection{Prior work}

Lages and Lopes proved the conjecture for profiles with three, four or
five cycles \cite[Theorem~1.7]{LL}.
Watanabe proved it for profiles $\{1,\ell_1,\ell_2\}$
\cite[Proposition~2.2]{W} and, using a list of connected quandles,
``In the case where the order of $Q$ is less than or equal to $47$''
\cite[Proposition~3.1]{W}. Kajiwara and Nakayama proved the affine case
first. Their summary reads ``We prove that the affine case of a conjecture
by C. Hayashi [\dots] that any connected finite quandle has a large orbit''
(zbMATH 1398.20081, for \cite{KN}), and Kayacan writes that ``Kajiwara and
Nakayama's result shows that Hayashi's Conjecture holds for connected
affine quandles'' \cite[\S1]{K}. Rehman's thesis also states that ``The
Hayashi's conjecture is true for any connected affine quandle''
\cite[Observations~1.4.7(4)]{Rt}, as does \cite[\S4, the second Case~4]{R}.
Both use a theorem of Hor\v{o}sevski\u{\i}~\cite{Ho}, which Bors states as
``every automorphism $\alpha$ of a finite nilpotent group has a cycle whose
length coincides with $\operatorname{ord}(\alpha)$'' \cite[Abstract]{B}.
Filipi records that a review of \cite{KN} pointed out that
Hor\v{o}sevski\u{\i}'s results ``also contribute'' \cite[\S1]{F}. Bors also
showed that ``the least order of a group where this property is violated
for an appropriate automorphism is $120$'' \cite[Abstract]{B}, so outside
nilpotent groups an automorphism can have no cycle as long as its order;
our right translations are of this kind. Kayacan proved the conjecture for
racks on which the inner automorphism group acts primitively
\cite[Theorem~6]{K}.

For conjugation quandles, Filipi records that ``this conjecture still
remains open'' \cite[\S1]{F}, and writes after Conjecture~4.1 that it was
``verified for all finite nilpotent groups, for the groups $S_n$, $A_n$,
$D_{2n}$, for all group of order $\le500$, and for all simple groups of
order $\le3.500.000$'' \cite[\S4]{F}. Kaur, Sharma and Singh proved that
for ``all finite irreducible Coxeter groups $G$ with a conjugation closed
subset $C$, the conjugation quandle $\mathrm{Conj}(G,C)$ satisfies the
Hayashi property'' \cite[Abstract]{KSS}. Spaggiari's thesis, in a chapter
based on a joint project with Bonatto, shows that the conjecture holds for
all finite connected quandles if it holds for the faithful ones
\cite[Theorem~6.4.1]{S}. The quandles of Theorem~\ref{thm:main} are
faithful, so they also refute that reduced form, and they generate groups
of order greater than $500$, beyond the range Filipi checked. The joint
paper of Spaggiari and Bonatto, announced in \cite{S} as work in progress,
had not been posted by November 5, 2026.

\subsection{What is new}

The counterexamples, the family and the bound of
Proposition~\ref{prop:min} are new. A right translation $R$ fails
Hayashi's condition exactly when no cycle of $R$ has length equal to its
order. By the theorem of Hor\v{o}sevski\u{\i} this cannot happen for an
automorphism of a finite nilpotent group, which is how the affine case is
settled. Our quandles live in groups of affine maps that are not
nilpotent. The step that needed a new idea is to make every cycle of $R$
too short while keeping the quandle connected. The cycle lengths of $R$ are
governed by vectors $d_b$ attached to the linear parts. Connectedness comes
from the translations by vectors fixed by $P$: they commute with $c_0$, so
adding them to the group generated by $C$ does not change the class, and
with them that group fills the whole affine group (Lemma~\ref{lem:conn}).
It remains open whether some connected quandle of order between $48$ and
$1889$ refutes the conjecture; by Proposition~\ref{prop:min} it would not
come from our construction. The proofs below are complete by hand;
Section~\ref{sec:verification} describes the computer checks of
Theorem~\ref{thm:main} and Corollary~\ref{cor:filipi}.

\section{Proof overview}

Let $B$ be a direct sum of vector spaces $V_q$ over $\F_q$ for distinct
primes $q$, let $A\cong\F_3^2$ act on $B$, and let $P$ be an involution
with $PaP=a^{-1}$ for $a\in A$. The quandle is the class $C$ of the affine
map $c_0\colon w\mapsto P(w+g_0)$ under conjugation by the affine maps with
linear part in $A$. Every element of $C$ has linear part $bP$ for some
$b\in A$; the \emph{fiber} of $b$ is the set of elements of $C$ with linear
part $bP$. The map $R=R_{c_0}$ swaps the fibers of $b$ and $b^{-1}$. On the
fiber of $b\ne1$, the square $R^2$ is a translation by a vector $d_b$, so
the cycles there have length $2\operatorname{ord}(d_b)$, where
$\operatorname{ord}(d_b)$ is the product of the primes $q$ at which $d_b$
has a nonzero component.

We choose the data so that $d_b$ has no component at $q=2,3,5,7$
respectively when $b$ lies on the first, second, third or fourth of the
four lines through $0$ (subgroups of order $3$) of $\F_3^2$. Every $b\ne1$
lies on exactly one line, so the cycle lengths outside the fiber $b=1$ are
$210$, $140$, $84$ and $60$, and none equals $\operatorname{ord}(R)=420$.
Connectedness and faithfulness follow from a short computation in the
group (Lemmas~\ref{lem:conn} and~\ref{lem:faith}).
Proposition~\ref{prop:min} rests on the fact that a group is never a union
of three proper subgroups of index at least~$3$.

\section{The construction}\label{sec:setting}

\paragraph{Data.} Let $B=\bigoplus_q V_q$, where $q$ runs over a finite set
of distinct primes and $V_q$ is a finite-dimensional $\F_q$-vector space;
write $v_q$ for the $V_q$-component of $v\in B$. All linear maps below map
each $V_q$ to itself. Let $A$ be a group of such invertible linear maps, of
odd order, let $P$ be such a linear map with $P^2=1$ and $PaP=a^{-1}$ for
all $a\in A$, and let $g_0\in B$.

The group $A$ is abelian, because $a\mapsto PaP=a^{-1}$ is a homomorphism.
Each $b\in A$ has a unique square root $a=b^{(|A|+1)/2}$, and then
$aPa^{-1}=a(Pa^{-1}P)P=a^2P=bP$, so $(bP)^2=1$.

Write $(Y,u)$ for the affine map $w\mapsto Yw+u$ of $B$, so that
$(Y,u)(Y',v)=(YY',Yv+u)$ and $(Y,u)^{-1}=(Y^{-1},-Y^{-1}u)$. Let
\[
G=\{(a,v):a\in A,\ v\in B\},\qquad c_0=(P,Pg_0),\qquad
C=\{gc_0g^{-1}:g\in G\},
\]
with $x*y=yxy^{-1}$ on $C$. For a subgroup or element $X$ of $A\langle
P\rangle$ and a subspace $V$, let $C_V(X)$ be the vectors of $V$ fixed by
$X$, and let $[V,A]$ be the span of the vectors $(1-a)v$. Put
\[
h=(1+P)g_0,\qquad d_b=(1-b)h\quad(b\in A),\qquad U_b=(1-bP)B.
\]
Multiplying out gives two formulas, used throughout. For $g=(a,v)\in G$
and $b\in A$,
\begin{align}
gc_0g^{-1}&=(a^2P,\ (1-a^2P)v+aPg_0),\label{eq:E1}\\
c_0(a,u)c_0^{-1}&=(a^{-1},\ Pu+P(1-a)g_0),\qquad
c_0(bP,u)c_0^{-1}=(b^{-1}P,\ Pu+Pg_0-b^{-1}g_0).\label{eq:E2}
\end{align}
For~\eqref{eq:E2}, $c_0^{-1}=(P,-g_0)$ and $c_0(Y,u)c_0^{-1}=(PYP,\
-PYg_0+Pu+Pg_0)$; put $Y=a$, or $Y=bP$ and use $Pb=b^{-1}P$.

\begin{lemma}[closure]\label{lem:closure}
Conjugation by $c_0$ maps $G$ onto $G$, and $C$ is invariant under
conjugation by every element of $\hat G=G\langle c_0\rangle\supseteq C$.
So $C$ is a conjugation quandle, and all its right translations are
conjugate, as permutations of $C$, to $R=R_{c_0}$.
\end{lemma}

\begin{proof}
By~\eqref{eq:E2}, $c_0Gc_0^{-1}\subseteq G$, with equality since $G$ is
finite. Hence $c_0(gc_0g^{-1})c_0^{-1}=g'c_0g'^{-1}$ with
$g'=c_0gc_0^{-1}\in G$, and $C$ is invariant under $G$ and $c_0$, hence
also under $c_0^{-1}$, since conjugation by $c_0$ is injective and $C$ is
finite. For $y=kc_0k^{-1}$ with $k\in G$ we have $R_y=\varphi
R\varphi^{-1}$, where $\varphi$ is conjugation by $k$.
\end{proof}

\begin{lemma}[the elements]\label{lem:elements}
$C=\{(bP,u): b\in A,\ u\in U_b+aPg_0\}$, where $a$ is the square root of
$b$, and $|C|=|A|\,|U_1|=|A|\,|B|/|C_B(P)|$.
\end{lemma}

\begin{proof}
In~\eqref{eq:E1}, $b=a^2$ runs over $A$ and $v$ over $B$. Since
$bP=aPa^{-1}$, we have $U_b=aU_1$, and $|U_1|=|B|/|\ker(1-P)|$.
\end{proof}

\begin{lemma}[the cycles of $R$]\label{lem:cycles}
Let $x=(bP,u)\in C$.
\begin{enumerate}
\item $R(x)=(b^{-1}P,\ Pu+Pg_0-b^{-1}g_0)$ and $R^2(x)=(bP,\ u+d_b)$.
\item If $b\ne1$, the cycle of $x$ has length $2\prod\{q:(d_b)_q\ne0\}$.
\item If $b=1$, write $u=Pg_0+w$ with $w\in U_1$. Then $R(x)=(P,Pg_0-w)$,
so $R$ fixes the $|U_1\cap V_2|$ points with $2w=0$ (one point if $2$ is
not among the primes) and swaps the others in pairs.
\end{enumerate}
\end{lemma}

\begin{proof}
(1) The first formula is~\eqref{eq:E2}. Applying it twice, with
$Pb^{-1}=bP$, gives $u\mapsto u+g_0-bPg_0+Pg_0-bg_0=u+(1-b)(1+P)g_0$.
(2) Since $|A|$ is odd, $b^{-1}\ne b$; odd powers of $R$ move $x$ to the
fiber of $b^{-1}$, and $R^{2m}(x)=(bP,u+md_b)$. The additive order of $d_b$
is the product of the primes at which it is nonzero.
(3) Since $P(1-P)=P-1$, we have $Pw=-w$ for $w\in U_1$, so
$Pu+Pg_0-g_0=Pg_0-w$. The elements of order at most $2$ in $B$ are those of
$V_2$.
\end{proof}

\begin{lemma}[connectedness]\label{lem:conn}
If $2$ is not among the primes, or some $a\in A$ fixes no nonzero vector of
$V_2$, then $C$ is connected.
\end{lemma}

\begin{proof}
First, $[B,A]+(1-P)B+C_B(P)=B$. For odd $q$ every $v\in V_q$ is
$\frac12(1-P)v+\frac12(1+P)v$, and $(1+P)v$ is fixed by $P$; on $V_2$, if
$a$ fixes no nonzero vector then $1-a$ is invertible and $[V_2,A]=V_2$.

Let $\alpha(g)=c_0gc_0^{-1}$, an automorphism of $G$, and let $N$ be the
subgroup generated by the elements $g\alpha(g)^{-1}=(gc_0g^{-1})c_0^{-1}$,
$g\in G$. Then $N\le H=\langle C\rangle$, and $N$ is normal in $G$ because
$k\,g\alpha(g)^{-1}k^{-1}=(kg)\alpha(kg)^{-1}\cdot(k\alpha(k)^{-1})^{-1}$.
The linear part of $(a,v)\alpha((a,v))^{-1}$ is $a^2$, so the linear parts
of $N$ are all of $A$. For $n=(a,r)\in N$ and $t\in B$, the commutator
$(1,t)n(1,t)^{-1}n^{-1}=(1,(1-a)t)$ lies in $N$, so $N$ contains the
translations by $[B,A]$. It also contains $(1,v)\alpha((1,v))^{-1}=
(1,(1-P)v)$. The translations by $C_B(P)$ form a group $T$ that commutes
with $c_0$. By the first paragraph $NT$ contains every translation, and it
has every linear part, so $NT=G$. Hence every $gc_0g^{-1}$ equals
$nc_0n^{-1}$ for some $n\in N\subseteq H$, and conjugation by an element of
$H$ is a product of maps $R_y^{\pm1}$ with $y\in C$.
\end{proof}

\begin{lemma}[faithfulness]\label{lem:faith}
If $C_{V_2}(A)\cap C_{V_2}(P)=0$, or $2$ is not among the primes, then $C$
is faithful.
\end{lemma}

\begin{proof}
Suppose $R_x=R_y$ with $x=(bP,u)$ and $y=(b'P,u')$. Then $z=y^{-1}x$
commutes with every element of $C$, and $z\in G$ has linear part $b'b^{-1}$.
As $z$ commutes with $c_0$, \eqref{eq:E2} shows its linear part equals its
inverse, so it is $1$; thus $b=b'$ and $z=(1,s)$. From $x=yz=(bP,u'+bPs)$
we get $bPs=u-u'\in U_b$. Conjugating $(1,s)$ by $(b''P,u'')\in C$ gives
$(1,b''Ps)$, so $b''Ps=s$ for all $b''\in A$; hence $s\in C_B(P)\cap
C_B(A)$ and $s=bPs\in U_b$. On $U_b$ the map $bP$ acts as $-1$, so $s=-s$
and $s_q=0$ for odd $q$, while $s_2\in C_{V_2}(A)\cap C_{V_2}(P)=0$.
Hence $s=0$, so $z=1$ and $x=y$.
\end{proof}

\section{The counterexamples}\label{sec:examples}

\begin{proof}[Proof of Theorem~\ref{thm:main}]
Take $A=\langle a_1,a_2\rangle\cong\F_3^2$ acting on
$B=V_2\oplus V_3\oplus V_5\oplus V_7$, with $P$ and $g_0$ as in
Appendix~\ref{app:data}, where the hypotheses of Section~\ref{sec:setting}
are checked by hand. For $q\in\{2,3,5,7\}$ let $K_q=\{b\in A:(d_b)_q=0\}$,
a subgroup of $A$ since it is the stabilizer of $h_q$. Appendix~\ref{app:data}
shows that $K_2,K_3,K_5,K_7$ are the four distinct lines of $\F_3^2$. Every
$b\ne1$ lies on exactly one line, so by Lemma~\ref{lem:cycles}(2) its fiber
consists of cycles of length $2\cdot210/q$, where $q$ is the prime of its
line: $210$, $140$, $84$ or $60$ for $q=2,3,5,7$. Lemma~\ref{lem:cycles}(3)
gives the cycles of lengths $1$ and $2$ on the fiber $b=1$. The counts in
the profiles follow from Lemma~\ref{lem:elements} and the sizes computed
in the appendix: in the quandle of order $1890$ each fiber has $210$
points, and in that of order $5670$ each has $630$. Connectedness and
faithfulness follow from Lemmas~\ref{lem:conn} and~\ref{lem:faith},
because $a_1$ fixes no nonzero vector of $V_2$. By Lemmas~\ref{lem:closure}
and~\ref{lem:conn}, $C$ is a single conjugacy class of $\langle C\rangle$.
\end{proof}

\begin{proof}[Proof of Corollary~\ref{cor:filipi}]
By Lemmas~\ref{lem:closure} and~\ref{lem:conn}, $C$ is a finite conjugacy
class of $H=\langle C\rangle$. Fix $c,z\in C$, let $\ell$ be the length of
the cycle of $R_c$ through $z$, and let $m$ be the order of $R_c$, the
least common multiple of its cycle lengths. Since $c^kzc^{-k}=R_c^k(z)$,
the power $c^\ell$ commutes with $z$. It is central in $H$ if and only if
it commutes with every element of $C$, that is, if and only if $m\mid\ell$.
Every right translation has the profile of Theorem~\ref{thm:main}, so
$m=420$ and $\ell\le210$. Thus $c^\ell$ commutes with $z$ without being
central, and no $z$ satisfies the conclusion of the conjecture.
\end{proof}

The same argument shows that every quandle of Theorem~\ref{thm:family}
refutes Filipi's conjecture, since there the order $2\prod_jq_j$ of $R$
exceeds every $\ell_j$. The argument is the equivalence of
\cite[Theorem~4.1]{F}, applied to a single class.

\section{The family and latin examples}

\begin{proof}[Proof of Theorem~\ref{thm:family}]
Let $A=\F_p^2$, with its $p+1$ subgroups of order $p$ (the lines)
$L_0,\dots,L_p$, and choose surjective homomorphisms $\chi_j\colon
A\to\mathbb{Z}/p$ with kernel $L_j$. Let $\Pi=1+x+\dots+x^{p-1}$. For each
$j$ take the block $V_{q_j}$ over $\F_{q_j}$, together with the actions of
$a\in A$ and of $P$ and the component of $g_0$, as follows:
\begin{itemize}
\item if $q_j\ne p$: $V_{q_j}=\F_{q_j}[x]/(\Pi)$, $a$ acts as
multiplication by $x^{\chi_j(a)}$, $Pf(x)=f(x^{-1})$, and $g_0$ has
component $x$; $P$ is well defined because $\Pi(x^{-1})=x^{1-p}\Pi(x)$;
\item if $q_j=p$: $V_p=\F_p^2$, $a$ acts as $J^{\chi_j(a)}$ with
$J=\left(\begin{smallmatrix}1&1\\0&1\end{smallmatrix}\right)$,
$P=\operatorname{diag}(-1,1)$, and $g_0$ has component $(0,1)$.
\end{itemize}
In both cases $P^2=1$ and $PaP=a^{-1}$, and $A$ acts faithfully because
$\bigcap_jL_j=1$.

If $q_j\ne p$, then $h_{q_j}=x+x^{-1}=x+x^{p-1}$ is nonzero in $V_{q_j}$,
since for $p\ge3$ it has two nonzero coefficients and is not a multiple of
$\Pi$, which has $p$. For $p\nmid k$, the polynomial $1-x^k$ is a unit
modulo $\Pi$, since $\gcd(x^k-1,x^p-1)=x-1$ and $\Pi(1)=p\ne0$ in
$\F_{q_j}$. So $(d_b)_{q_j}=(1-b)h_{q_j}$ is nonzero exactly when $b\notin
L_j$. If $q_j=p$, then $h_p=(0,2)$ and $(1-J^k)h_p=(-2k,0)$, again nonzero
exactly when $p\nmid k$. Hence for $b\ne1$ on $L_j$, $d_b$ is nonzero
exactly at the primes $q_i$ with $i\ne j$, and Lemma~\ref{lem:cycles} gives
cycles of length $\ell_j$. The remaining lengths are $1$ and $2$.

If some $q_j=2$, then $V_2=\F_2[x]/(\Pi)$ and every $a\notin L_j$ acts on
it without nonzero fixed vectors, since $1-x^k$ is a unit; so
Lemmas~\ref{lem:conn} and~\ref{lem:faith} apply, as they do when no $q_j$
is $2$. For the order, $|C|=p^2\prod_j|(1-P)V_{q_j}|$ by
Lemma~\ref{lem:elements}. If $q_j=p$, then $1-P=\operatorname{diag}(2,0)$
has rank $1$. If $q_j\ne p$, the image of $1-P$ on
$\F_{q_j}[x]/(x^p-1)$ is spanned by the $(p-1)/2$ elements $x^k-x^{-k}$,
$1\le k\le (p-1)/2$, which have disjoint supports; their span meets the
kernel $\F_{q_j}\Pi$ of the map to $V_{q_j}$ only in $0$, since its elements
have constant term $0$. So $(1-P)V_{q_j}$ has dimension $(p-1)/2$.

For $j\ne k$, $\ell_j$ does not divide $\ell_k$: if $q_k$ is odd, $q_k$
divides $\ell_j$ but not $\ell_k$; if $q_k=2$, then $4$ divides $\ell_j$
but not $\ell_k$. So the $\ell_j$ are distinct, and none of the other $p$
divides the longest.

Suppose finally that all $q_j$ are odd, and let $y=(bP,u)\in C$. Then
$yc_0y^{-1}=(b^2P,(1-b^2P)u+bg_0)$. If $y$ and $y'=(b'P,u')$ have the same
image, then $b^2=b'^2$, so $b=b'$, and $(1-b^2P)v=0$ for $v=u-u'\in U_b$.
On $U_b$ we have $bPv=-v$, so $(1-b^2P)v=(1+b)v$. If $(1+b)v=0$ then
$v=b^{|A|}v=-v$, so $2v=0$ and $v=0$. Thus the left translation $y\mapsto
yc_0y^{-1}$ is injective, hence bijective, on $C$. Every element of $C$ is
$kc_0k^{-1}$ with $k\in G$, and its left translation is that of $c_0$
conjugated by $k$. So $C$ is latin.
\end{proof}


\section{Computational verification}\label{sec:verification}

Theorem~\ref{thm:main} and Corollary~\ref{cor:filipi} are also checked by
computer for both quandles, from the data of Appendix~\ref{app:data}. The
data are those of Section~\ref{sec:examples}, not the bases of
Theorem~\ref{thm:family}, whose latin examples rest on the proof above.
Four checks are used; none of their programs imports another. The first
three start from the input data, and the fourth reads the table built by
the second.
\begin{itemize}
\item The first, in Python with exact arithmetic over $\F_q$, checks that
$c_0$ normalizes $G$, computes the class $C$ and the cycle type of
$R_{c_0}$, and shows that the right translations by six elements of $C$
already act transitively. It forms no operation table.
\item The second is a Python program that rebuilds $C$ and its full
operation table, rejecting any product outside $C$, followed by a C program
that checks the quandle axioms on the table: idempotence, bijectivity of
every $R_y$, and right self-distributivity over all $n^3$ triples. It then
checks connectedness, that all $n$ right translations have the stated cycle
type, and faithfulness.
\item The third checks Filipi's statement in the group: it computes the
powers of $c_0$, finds that only $c_0^0$ is central in $\langle C\rangle$,
and for every $z\in C$ finds a $k\le210$ with $c_0^k$ commuting with $z$.
\item The fourth reads the table and shows that, for every $c\in C$, no
cycle of $R_c$ has length equal to its order.
\end{itemize}
Each check rejects an altered input aimed at it. With the $V_3$-component of
$g_0$ set to zero, the profile satisfies the conjecture, as
Lemma~\ref{lem:cycles} predicts (for order $5670$ it becomes
$1^2\,2^{314}\,20^{63}\,28^{45}\,70^{18}\,140^{9}$); the first, the C
program and the third and fourth checks reject it. With an identity block
on $V_7$ in $P$, $c_0$ no longer normalizes $G$, and the first check and the
table builder reject it. With one table entry changed, the C program
rejects the table. On an ordinary computer all checks take about three
minutes, and the complete distributivity check for order $5670$ about
seven minutes more. The programs and commands are distributed with this
paper.

\begin{appendices}
\section{The data and their hand verification}\label{app:data}

Vectors are columns and matrices are given by rows. The quandle of order
$5670$ uses $V_2=\F_2^2$, $V_3=\F_3^3$, $V_5=\F_5^2$, $V_7=\F_7^2$ and
\[
\begin{array}{c|cccc}
 & a_1 & a_2 & P & g_0\\\hline
\F_2^2 & E=\begin{pmatrix}0&1\\1&1\end{pmatrix} & I & \begin{pmatrix}1&1\\0&1\end{pmatrix} & (0,1)\\[8pt]
\F_3^3 & I & \sigma=\begin{pmatrix}0&0&1\\1&0&0\\0&1&0\end{pmatrix} & -\begin{pmatrix}1&0&0\\0&0&1\\0&1&0\end{pmatrix} & (0,1,0)\\[12pt]
\F_5^2 & M=\begin{pmatrix}0&4\\1&4\end{pmatrix} & M & \begin{pmatrix}1&4\\0&4\end{pmatrix} & (1,0)\\[8pt]
\F_7^2 & \operatorname{diag}(2,4) & \operatorname{diag}(4,2) & \begin{pmatrix}0&1\\1&0\end{pmatrix} & (1,0)
\end{array}
\]
The quandle of order $1890$ changes only the $\F_3$-block, to $a_1=I$,
$a_2=J\oplus1$ with $J=\left(\begin{smallmatrix}1&1\\0&1\end{smallmatrix}\right)$,
$P=\operatorname{diag}(2,1,1)$ and $g_0=(0,1,0)$. Its third $\F_3$-coordinate
is inert ($A$ and $P$ fix it and $g_0$ vanishes there); it keeps the block
sizes of the order-$5670$ data.

\paragraph{$A\cong\F_3^2$.} In each block $a_1$ and $a_2$ commute (one is
$I$, or they are equal, or both are diagonal). Each has order $3$:
$E^3=I$ over $\F_2$, $\sigma$ is a $3$-cycle of coordinates, $J^3=I$ over
$\F_3$, $M^3=I$ over $\F_5$, and $2^3=4^3=1$ in $\F_7$. If $a_1^ia_2^k=1$,
the $\F_2$-block gives $3\mid i$ and the $\F_3$-block gives $3\mid k$. So
$|A|=9$.

\paragraph{$P^2=1$ and $Pa_iP=a_i^{-1}$.} Over $\F_2$: $P^2=I$ and
$PEP=\left(\begin{smallmatrix}1&1\\1&0\end{smallmatrix}\right)=E^2$. Over
$\F_3$: conjugating the $3$-cycle $\sigma$ by the transposition of the last
two coordinates gives $\sigma^{-1}$; for the order-$1890$ data,
$\operatorname{diag}(-1,1)J\operatorname{diag}(-1,1)=J^{-1}$. Over $\F_5$:
$P^2=I$ and $PMP=\left(\begin{smallmatrix}4&1\\4&0\end{smallmatrix}\right)=
M^{-1}$. Over $\F_7$: swapping coordinates turns $\operatorname{diag}(2,4)$
into $\operatorname{diag}(4,2)=\operatorname{diag}(2,4)^{-1}$.

\paragraph{The lines $K_q$.} The components of $h=(1+P)g_0$ are
$h_2=(1,0)$, $h_3=(0,1,2)$ (order $5670$) or $(0,2,0)$ (order $1890$),
$h_5=(2,0)$ and $h_7=(1,1)$. Let $b=a_1^ia_2^k\ne1$.
\begin{itemize}
\item $q=2$: $b$ acts as $E^i$, which for $3\nmid i$ fixes no nonzero
vector, since $x^2+x+1$ has no root in $\F_2$. So $K_2=\langle a_2\rangle$.
\item $q=3$: for order $5670$, $b$ acts as $\sigma^k$, whose fixed vectors
for $3\nmid k$ are the constant vectors, and $h_3$ is not constant; for
order $1890$, $(1-b)h_3=(k,0,0)$. So $K_3=\langle a_1\rangle$.
\item $q=5$: $b$ acts as $M^{i+k}$, which for $i+k\not\equiv0$ fixes no
nonzero vector, since $x^2+x+1$ has no root in $\F_5$. So
$K_5=\langle a_1a_2^2\rangle$.
\item $q=7$: $b$ acts as $\operatorname{diag}(2^{i+2k},2^{2i+k})$, and both
coordinates of $h_7$ are nonzero, so $(d_b)_7=0$ needs $3\mid i+2k$ and
$3\mid 2i+k$, since $2^j=1$ exactly when $3\mid j$. Both congruences say
$i\equiv k\pmod3$, so $K_7=\langle a_1a_2\rangle$.
\end{itemize}

\paragraph{Counting.} The fixed spaces of $P$ are $\{(x,0)\}$ in $\F_2^2$,
$\{(0,2z,z)\}$ in $\F_3^3$ (order $5670$) or $\{(0,y,z)\}$ (order $1890$),
$\{(x,0)\}$ in $\F_5^2$ and the diagonal in $\F_7^2$. So $|C_B(P)|$ is
$2\cdot3\cdot5\cdot7=210$ or $2\cdot9\cdot5\cdot7=630$, while
$|B|=4\cdot27\cdot25\cdot49=132300$. Hence $|U_1|$ is $630$ or $210$, and
$|C|=9|U_1|$ is $5670$ or $1890$. In both cases $U_1\cap V_2=(1-P)V_2$ has
$2$ elements, so the fiber $b=1$ contributes $1^2\,2^{(|U_1|-2)/2}$. Each
line contributes two fibers of $|U_1|$ points, divided into cycles of the
length of that line.

\section{Proof of Proposition~\ref{prop:min}}\label{app:min}

\begin{proof}[Proof of Proposition~\ref{prop:min}]
Let $S(b)$ be the set of primes $q$ with $(d_b)_q\ne0$, and $S=\bigcup_b
S(b)$. By Lemma~\ref{lem:cycles} the cycle lengths of $R$ are $1$, possibly
$2$, and $2\prod S(b)$ for $b\ne1$. If $S(b^*)=S$ for some $b^*\ne1$, every
length divides $2\prod S$, which is itself a length. So if some length
fails, every $b\in A$ lies in one of the proper subgroups
$K_q=\{b:(d_b)_q=0\}$ with $q\in S$. A cyclic group is not a union of
proper subgroups, and in a group of odd order a proper subgroup has at most
$|A|/3$ elements, so three of them cover at most $|A|-2$ elements. Hence
$|S|\ge4$ and $A$ is non-cyclic of odd order, so $|A|\ge9$ and $\prod S\ge
2\cdot3\cdot5\cdot7$.

For $q\in S$, pick $b$ with $q\in S(b)$. Both $u$ and $u+d_b$ lie in
$U_b+aPg_0$ by Lemma~\ref{lem:cycles}(1), so $(d_b)_q$ is a nonzero element
of $(1-bP)V_q$, which has the size of $(1-P)V_q$. So $|U_1|\ge\prod S\ge210$
and $|C|=|A|\,|U_1|\ge1890$.

In a latin quandle $R_y$ fixes only $y$, since $x*y=x=x*x$ forces $x=y$. By
Lemma~\ref{lem:cycles}(3), $U_1\cap V_2=(1-P)V_2$ is then $0$, so $2\notin
S$, $\prod S\ge3\cdot5\cdot7\cdot11$ and $|C|\ge9\cdot1155=10395$.
Theorem~\ref{thm:main} and the case $p=3$, $\{q_j\}=\{3,5,7,11\}$ of
Theorem~\ref{thm:family} attain the bounds.
\end{proof}

\end{appendices}

\begin{thebibliography}{99}
\bibitem{B} A. Bors, \emph{On finite groups where the order of every
automorphism is a cycle length}, arXiv:1412.8418v2 (2015).
\bibitem{F} F. Filipi, \emph{Hayashi property for conjugation quandles},
arXiv:2503.22377v1 (2025).
\bibitem{Ha} C. Hayashi, \emph{Canonical forms for operation tables of
finite connected quandles}, Comm. Algebra 41 (2013), no.~9, 3340--3349,
\href{https://doi.org/10.1080/00927872.2012.685532}{doi:10.1080/00927872.2012.685532};
arXiv:1110.1849v1 (2011).
\bibitem{Ho} M. V. Hor\v{o}sevski\u{\i}, \emph{On automorphisms of finite
groups}, Math. USSR-Sb. 22 (1974), no.~4, 584--594,
\href{https://doi.org/10.1070/SM1974v022n04ABEH001707}{doi:10.1070/SM1974v022n04ABEH001707}.
\bibitem{KN} T. Kajiwara and C. Nakayama, \emph{A large orbit in a finite
affine quandle}, Yokohama Math. J. 62 (2016), 25--29; summary in zbMATH
1398.20081.
\bibitem{KSS} D. Kaur, U. B. Sharma and P. Singh, \emph{Hayashi property for
conjugation quandles over finite Coxeter groups}, arXiv:2607.03922v1 (2026).
\bibitem{K} S. Kayacan, \emph{On a conjecture about profiles of finite
connected racks}, Erzincan University Journal of Science and Technology 18
(2025), no.~3, 735--744,
\href{https://doi.org/10.18185/erzifbed.1591413}{doi:10.18185/erzifbed.1591413};
arXiv:2006.10327v3.
\bibitem{LL} A. Lages and P. Lopes, \emph{On a conjecture by Hayashi on
finite connected quandles}, Ars Math. Contemp. 26 (2026), no.~4, Paper
No.~P4.05,
\href{https://doi.org/10.26493/1855-3974.3529.4e5}{doi:10.26493/1855-3974.3529.4e5}.
Statement numbers are those of the published version; in arXiv:2405.11660v2
they are Conjecture~1.1, Theorem~1.3 and Corollary~2.1.
\bibitem{R} N. Rehman, \emph{On the cycle structure of finite racks and
quandles}, arXiv:1912.05115v1 (2019).
\bibitem{Rt} N. Rehman, \emph{Quandles and Hurwitz orbits}, PhD thesis,
Philipps-Universit\"at Marburg, 2016,
\url{https://d-nb.info/1114394920/34}.
\bibitem{S} F. Spaggiari, \emph{Non-associative algebraic structures}, PhD
thesis, Charles University, Prague, 2025.
\bibitem{W} T. Watanabe, \emph{On the structure of the profile of finite
connected quandles}, Math. J. Okayama Univ. 61 (2019), 85--98,
\href{https://doi.org/10.18926/mjou/56328}{doi:10.18926/mjou/56328},
\url{https://www.math.okayama-u.ac.jp/mjou/mjou61/_05_Watanabe.pdf}.
\end{thebibliography}
\end{document}
